\section{Experiment Evaluation and Results}
\label{sec:evaluation}

This section evaluates the DeSFAM framework in terms of syscall-level enforcement accuracy,
performance overhead, reduction rate, and scalability in containerized environments.
Experiments were conducted using real-world workloads and attack scenarios to assess DeSFAM's
effectiveness and runtime efficiency. Metrics are compared against leading frameworks under
controlled conditions to ensure fairness and reproducibility. Results highlight DeSFAM's
security impact and operational feasibility in practical deployments.

\subsection{Experimental Setup}
\label{subsec:experimental_setup}

\subsubsection{Infrastructure}
\label{subsubsec:infrastructure}

Experiments were conducted on a high-performance bare-metal server running Ubuntu~24.04~LTS
equipped with a 16-core CPU, 64~GB RAM, NVIDIA A100 80~GB GPU, and a 512~GB NVMe SSD. Docker
26.0 with containerd~1.7 was used for container lifecycle management. The eBPF environment was
configured with libbpf~1.3 and bpftrace~0.21. Model training was performed using PyTorch~2.3
with CUDA acceleration.

\subsubsection{Container Workloads}
\label{subsubsec:container_workloads}

We evaluated DeSFAM on six diverse containers: Nginx (web server), MySQL~8.0 (relational
database), Redis~7.2 (key-value store), Node.js~20 (JavaScript runtime), Python~3.12 (script
processing), and MongoDB~7.0 (document database)---spanning web serving, data storage, and
script execution with significantly different syscall profiles.

\subsubsection{Workload Generation}
\label{subsubsec:workload_generation}

Realistic traffic was generated using industry-standard tools: \texttt{wrk} and ApacheBench
(\texttt{ab}) for HTTP load on Nginx; \texttt{redis-benchmark} for Redis; and custom Python
scripts to simulate CRUD operations for MySQL and MongoDB. Each workload was run at three load
levels: idle, medium load (50\% capacity), and peak load (90\% capacity), in 10-minute
intervals.

\subsubsection{Attack Simulation}
\label{subsubsec:attack_simulation}

Three categories of real-world CVE attack scenarios were used:
\begin{itemize}
  \item \textbf{CVE-2021-4034 (Polkit pkexec)~\cite{cve_2021_4034}:} Local privilege escalation
        exploit via a crafted \texttt{execve} sequence.
  \item \textbf{CVE-2022-0847 (Dirty Pipe)~\cite{cve_2022_0847}:} Kernel memory overwrite
        exploit via \texttt{splice} and \texttt{pipe}.
  \item \textbf{CVE-2017-5123 (waitid exploit):} Kernel privilege escalation via \texttt{waitid}
        with an invalid pointer.
\end{itemize}

A total of 14 distinct CVE attack scenarios were executed across all container workloads.

\subsubsection{Baseline Comparisons}
\label{subsubsec:baselines}

DeSFAM was benchmarked against five baseline systems: Confine~\cite{ref31_confine},
Prof-gen~\cite{ref36_profgen}, NIMOS~\cite{ref38_nimos}, Optimus~\cite{ref27_optimus}, and
KubeRosy~\cite{ref29_kuberosy}. Each baseline was reimplemented following its original
configuration and integrated with Docker/Kubernetes for parity in runtime enforcement and
syscall profiling capabilities.

\subsection{Evaluation Metrics}
\label{subsec:metrics}

\subsubsection{Syscall Reduction Rate (SRR)}
SRR quantifies DeSFAM's ability to reduce kernel exposure by minimizing allowed syscalls:
\begin{equation}
  \text{SRR} = \left(1 - \frac{|\text{DeSFAM Access List}|}{|\text{Docker Default Access List}|}\right) \times 100\%
  \label{eq:srr}
\end{equation}

Docker's default allows 414 syscalls. DeSFAM's hybrid profiling (static, dynamic,
CVE-informed) yields container-specific access lists, achieving higher SRRs than static-only or
dynamic-only methods.

\subsubsection{Anomaly Detection Performance}
We assess SyscallAD's accuracy using:
\begin{itemize}
  \item Recall (TPR) -- Proportion of attacks correctly identified
  \item Precision -- Fraction of flagged sequences that are actual threats
  \item F1-Score, AUC, and Average Precision (AP) -- Robustness in imbalanced data
\end{itemize}
DeSFAM's VAE + iForest ensemble outperforms LSTM-only detection (e.g., NIMOS), with improved
precision and reduced false positives.

\paragraph{Reproduction on the DongTing dataset.}
We retrained SyscallAD's two components on the DongTing dataset
(18{,}966 sequences --- 6{,}850 benign, 12{,}116 malicious from 26 Linux
kernel releases; train=14{,}843 / val=1{,}812 / test=2{,}311). Features
follow the paper's 174-dimensional design: 60-syscall frequency vector,
40 discriminative binary flags, 8 distributional statistics, 40
bigram-transition rates, and a 26-way kernel-version one-hot. The
Isolation Forest is trained on the 5{,}487 normal sequences in the
training split; the VAE shares the same training set across 3 seeds with
the best-validation seed selected. Ensemble scores combine the two with
$\alpha = 0.7$ after a robust quantile-band normaliser (p1\,..\,p99 of
the training scores), which is required because raw VAE reconstruction
errors span six orders of magnitude on this dataset and a plain min--max
scaler would otherwise crush the VAE contribution. Table~\ref{tab:dongting_reproduction}
reports the metrics obtained.

\begin{table}[!t]
\centering
\caption{SyscallAD reproduction on the DongTing test split
(2{,}311 sequences, 70.7\% attack class).}
\label{tab:dongting_reproduction}
\begin{tabular}{lrrrrr}
\toprule
Model & AUC & AP & F1 & Precision & Recall \\
\midrule
Isolation Forest & 0.882 & 0.949 & 0.915 & 0.848 & 0.994 \\
VAE              & 0.960 & 0.970 & 0.963 & 0.937 & 0.990 \\
Ensemble ($\alpha{=}0.7$) & 0.921 & 0.960 & 0.932 & 0.876 & 0.996 \\
\bottomrule
\end{tabular}
\end{table}

The reproduction confirms the paper's headline finding that VAE +
Isolation Forest sequence-level anomaly detection is viable on
syscall-sequence data and reaches F1 above 0.92, but it also surfaces
one nuance: on the DongTing split the VAE component alone (F1 = 0.963,
AUC = 0.960) outperforms the $\alpha = 0.7$ ensemble (F1 = 0.932, AUC =
0.921) on every metric except recall. The ensemble's value on DongTing
is therefore in attack coverage (recall 99.6\%) rather than aggregate
F1; deployments that need lower false-positive rates may prefer to
deploy the VAE in isolation or to tune $\alpha$ downward.

\subsubsection{Attack Detection Latency}
Time (in $\mu$s) between malicious syscall execution and enforcement. DeSFAM's eBPF-LSM hooks
support inline blocking with sub-millisecond latency. Compared to Optimus and KubeRosy, which
rely on user-space components, DeSFAM reacts faster due to its kernel-level design.

\subsubsection{Enforcement Overhead}
We evaluate average CPU usage (\%), memory consumption (MB), and syscall processing latency
($\mu$s).

\subsubsection{Adaptability and Comparison Metrics}
To assess dynamic policy adaptation, we measured:
\begin{itemize}
  \item Rule update latency via SyscallAD or CVE triggers (in $\mu$s)
  \item Uptime impact (zero downtime observed)
\end{itemize}
None of the baseline systems---Confine, Prof-gen, NIMOS, Optimus, or KubeRosy---support
real-time policy self-hardening. For holistic comparison, we evaluate: access list size and
SRR, attack mitigation success rate, failure-free execution for benign workloads, and average
FPR and enforcement latency.

\subsection{Comparative Results}
\label{subsec:comparative_results}

We compare DeSFAM against five state-of-the-art syscall filtering frameworks---Confine,
Prof-gen, NIMOS, Optimus, and KubeRosy---using metrics such as syscall reduction rate (SRR),
enforcement stability, attack prevention, and runtime performance. The evaluation spans six
diverse container workloads: Nginx, MySQL, Redis, Node.js, Python, and MongoDB.

\subsubsection{Syscall Reduction Rate (SRR) Comparison}

A primary goal of DeSFAM is to minimize the kernel attack surface by reducing exposed syscalls.
Docker's default seccomp profile allows 414 syscalls, many of which are unnecessary and
potentially exploitable. DeSFAM employs hybrid reduction---static analysis, dynamic profiling,
and CVE-informed filtering---to generate container-specific access lists. The SRR is calculated
as:
\begin{equation}
  \text{SRR} = \left(1 - \frac{\text{Access List Size}}{414}\right) \times 100\%
  \label{eq:srr_comparative}
\end{equation}

Fig.~\ref{fig:srr_bar} and Fig.~\ref{fig:srr_rate} illustrate the improvement across all test
containers. Table~\ref{tab:srr} summarizes the results per framework.

\begin{figure}[htbp]
  \centering
  \includegraphics[width=\linewidth]{figure9}
  \caption{Syscall Access List Size and Reduction Rate Comparison.}
  \label{fig:srr_bar}
\end{figure}

\begin{figure}[htbp]
  \centering
  \includegraphics[width=\linewidth]{figure10}
  \caption{Syscall reduction rate comparison across frameworks.}
  \label{fig:srr_rate}
\end{figure}

\begin{table*}[tp]
  \centering
  \caption{Syscall Reduction Rate (SRR) Comparison.}
  \label{tab:srr}
  \includegraphics[width=\linewidth]{table4}
\end{table*}

DeSFAM achieves an average SRR of 78.60\%---the highest among all frameworks---while preserving
full container functionality. The average access list size is 89 syscalls, the smallest of all
compared systems.

\subsubsection{Functional Stability}
All six target containers maintained 100\% availability throughout all test workloads. No
functional regressions were observed, confirming that DeSFAM's access list synthesis balances
security and compatibility better than static or log-based access list generators.

\subsubsection{Attack Prevention}

DeSFAM's runtime blocking was evaluated against a set of real-world CVEs representing diverse
Linux kernel vulnerabilities. These CVEs cover multiple attack categories, including privilege
escalation, container escapes, namespace abuses, race conditions, and memory corruption exploits.
For safety and reproducibility, exploit behaviors were simulated by replaying the documented
syscall invocation sequences associated with each CVE, based on public disclosures and
proof-of-concept analyses.

Table~\ref{tab:cve_blocking} summarizes the blocking success across all frameworks. DeSFAM
successfully blocked all attack scenarios in real time using its SyscallAD sequence anomaly
detection and eBPF-LSM kernel-level enforcement. In contrast, baseline frameworks either failed
to detect these attacks or lacked full enforcement capability.

\begin{table*}[tp]
  \centering
  \caption{CVE Attack Blocking Success Rate.}
  \label{tab:cve_blocking}
  \includegraphics[width=\linewidth]{table5}
\end{table*}

\subsubsection{Baseline Comparison Metrics}

\begin{table*}[tp]
  \centering
  \caption{Baseline Comparison of Security and Performance Metrics.}
  \label{tab:baseline_comparison}
  \includegraphics[width=\linewidth]{table6}
\end{table*}

DeSFAM achieves the smallest access list size (89~syscalls), full CVE prevention (94\% of CVE
traces blocked), high benign stability (99\%), lowest FPR (1.6\%), and fastest enforcement
latency (0.540~$\mu$s).

\subsection{Runtime Performance}
\label{subsec:runtime_performance}

\subsubsection{CPU and Memory Overhead}

We evaluated resource usage across six representative containers with and without DeSFAM. As
shown in Fig.~\ref{fig:cpu_overhead} and Fig.~\ref{fig:mem_overhead}, DeSFAM incurs only a
1.74\% increase in average CPU utilization and reduces memory consumption by 6.1\%, primarily
due to blocked syscall paths and mitigation of exploit-induced memory bloat.

\begin{figure}[htbp]
  \centering
  \includegraphics[width=\linewidth]{figure11}
  \caption{CPU usage with and without DeSFAM. DeSFAM incurs minimal CPU overhead
           ($\approx$1.74\%) while enforcing real-time syscall policies.}
  \label{fig:cpu_overhead}
\end{figure}

\begin{figure}[htbp]
  \centering
  \includegraphics[width=\linewidth]{figure12}
  \caption{Memory usage with and without DeSFAM. Slight memory savings ($\approx$6.1\%)
           result from syscall suppression and exploit mitigation.}
  \label{fig:mem_overhead}
\end{figure}

\begin{table*}[tp]
  \centering
  \caption{CPU and Memory Usage With and Without DeSFAM.}
  \label{tab:cpu_memory}
  \includegraphics[width=\linewidth]{table7}
\end{table*}

Table~\ref{tab:overhead_summary} summarizes the aggregated performance overhead, confirming
DeSFAM's lightweight performance footprint with average memory savings of 50--55~MB.

\begin{table*}[tp]
  \centering
  \caption{Summary of Performance Overhead.}
  \label{tab:overhead_summary}
  \includegraphics[width=\linewidth]{table8}
\end{table*}

\subsubsection{Policy Update Latency}
To validate the responsiveness of DeSFAM's adaptive enforcement, we measured rule insertion and
propagation times in Phase~3. Upon anomaly detection, new rules were injected into the eBPF map
and applied in real time. Latency was measured using high-resolution kernel timestamps at both
the insertion and propagation stages.

Fig.~\ref{fig:policy_latency} presents the measured rule insertion and propagation latency
across six active containers. Table~\ref{tab:policy_latency} lists the per-container values.
The values were collected using high-resolution kernel timestamps before and after BPF map
updates and LSM hook propagation. All values remain consistently low (65--82~$\mu$s), confirming
the real-time capability of DeSFAM's in-kernel enforcement engine. All updates completed in
under 100~$\mu$s without interrupting service, demonstrating DeSFAM's capability for real-time,
zero-downtime syscall policy adaptation.

\begin{figure}[htbp]
  \centering
  \includegraphics[width=\linewidth]{figure13}
  \caption{DeSFAM Phase~3 policy update latency. All rules were inserted and enforced in under
           100~$\mu$s without container disruption.}
  \label{fig:policy_latency}
\end{figure}

\begin{table*}[tp]
  \centering
  \caption{Policy Update and Propagation Latency.}
  \label{tab:policy_latency}
  \includegraphics[width=\linewidth]{table9}
\end{table*}

\subsubsection{Syscall Latency Overhead}

We profiled syscall latency using \texttt{bpftrace} histograms for frequent operations such as
\texttt{read}, \texttt{write}, and \texttt{epoll\_wait}. Across all containers, DeSFAM
preserved sub-millisecond syscall latency, confirming its suitability for high-throughput
environments.

\begin{table}[htbp]
  \centering
  \caption{Syscall Latency Overhead (ns) for Common Operations.}
  \label{tab:syscall_latency}
  \begin{tabular}{lccc}
    \toprule
    \textbf{Syscall} & \textbf{Without DeSFAM} & \textbf{With DeSFAM} & \textbf{Delta} \\
    \midrule
    \texttt{read}       & 142\,ns & 218\,ns & +76\,ns \\
    \texttt{write}      & 156\,ns & 234\,ns & +78\,ns \\
    \texttt{epoll\_wait} & 89\,ns  & 148\,ns & +59\,ns \\
    \texttt{accept4}    & 201\,ns & 287\,ns & +86\,ns \\
    \bottomrule
  \end{tabular}
\end{table}

Table~\ref{tab:syscall_counts} summarizes the highest-frequency syscalls across each container
that completed within the 512--1000~$\mu$s range. Fig.~\ref{fig:syscall_latency_dist}
visualizes the aggregate counts per container in this sub-millisecond range. These results
demonstrate DeSFAM's ability to consistently maintain low syscall latency, even for frequent
operations like \texttt{read}, \texttt{write}, and \texttt{epoll\_wait} without introducing
measurable delay.

\begin{table*}[tp]
  \centering
  \caption{Syscall Counts in 512--1000\,$\mu$s Latency Range.}
  \label{tab:syscall_counts}
  \includegraphics[width=\linewidth]{table10}
\end{table*}

\begin{figure}[htbp]
  \centering
  \includegraphics[width=\linewidth]{figure14}
  \caption{Syscall latency distribution in the 512--1000~$\mu$s range. All syscall paths
           remained performant under DeSFAM.}
  \label{fig:syscall_latency_dist}
\end{figure}

DeSFAM preserves sub-millisecond syscall latency for all measured operations, with a maximum
delta of 86~ns for \texttt{accept4}. All values remain well under the 300~ns design target.

\subsubsection{Anomaly Detection Performance}

DeSFAM's SyscallAD model achieves strong classification metrics. The ensemble classifier reaches
96\% overall accuracy, with an F1-score of 0.92 on attack sequences. Fig.~\ref{fig:roc_pr}
shows ROC and PR curves, with AUC of 0.94 and AP of 0.87.

\begin{figure}[htbp]
  \centering
  \includegraphics[width=\linewidth]{figure15}
  \caption{SyscallAD model performance. Left: Precision-Recall Curve (AP~= 0.87). Right: ROC
           Curve (AUC~= 0.94).}
  \label{fig:roc_pr}
\end{figure}

\begin{table*}[tp]
  \centering
  \caption{SyscallAD Classification Metrics.}
  \label{tab:detection_performance}
  \includegraphics[width=\linewidth]{table11}
\end{table*}

SyscallAD achieves 96\% overall accuracy, F1-score of 0.92, and AUC of 0.94 on the DongTing
dataset. This represents a significant improvement over VAE-only (F1~= 0.84) or iForest-only
(F1~= 0.79) approaches, confirming the benefit of the ensemble fusion strategy.

\subsubsection{Feature Comparison with Existing Methods}
DeSFAM uniquely integrates real-time anomaly detection, adaptive policy enforcement, and syscall
minimization within a unified framework. Table~\ref{tab:feature_comparison} highlights its
advantage over existing frameworks across 12~technical dimensions.

\begin{table*}[tp]
  \centering
  \caption{Feature Comparison with Existing Frameworks.}
  \label{tab:feature_comparison}
  \includegraphics[width=\linewidth]{table12}
\end{table*}

\subsection{Ablation Study: Phase-Wise Contribution}
\label{subsec:ablation}

To understand the impact of each phase in DeSFAM, we conducted an ablation study disabling one
component at a time. As shown in Table~\ref{tab:ablation}, Phase~1 alone achieves 62.4\% syscall
reduction but blocks only 3 of 14 CVEs. Adding Phase~2 (SyscallAD) boosts anomaly detection,
raising the block count to 8 and reducing false positives significantly. Phase~3 provides
kernel-level enforcement, blocking 10 CVEs even without ML support. When all three phases are
active, DeSFAM blocks all 14 CVEs and reaches 81.6\% syscall reduction.

\begin{table}[htbp]
  \centering
  \caption{Ablation Study: Phase-Wise Effect on Security.}
  \label{tab:ablation}
  \includegraphics[width=\linewidth]{table13}
\end{table}

This study confirms that all phases are essential: Phase~1 minimizes exposure, Phase~2 enhances
detection, and Phase~3 ensures robust enforcement.
